\documentclass[11pt,a4paper]{report}

\usepackage[utf8]{inputenc}
\usepackage[T1]{fontenc}
\usepackage[french]{babel}
\usepackage[margin=2.2cm]{geometry}
\usepackage{xurl}
\usepackage{hyperref}
\usepackage{booktabs}
\usepackage{xcolor}
\usepackage{longtable}
\usepackage{enumitem}
\usepackage{fancyhdr}
\usepackage{titlesec}
\usepackage{array}
\usepackage{colortbl}
\usepackage{ragged2e}
\usepackage{listings}
\usepackage{tikz}

% ── Palette Copifi ──
\definecolor{copifiobsidian}{HTML}{050505}
\definecolor{copifiemerald}{HTML}{10B981}
\definecolor{copifisky}{HTML}{0EA5E9}
\definecolor{copifizinc500}{HTML}{71717A}
\definecolor{copificodebg}{HTML}{F4F4F5}
\definecolor{copificodeframe}{HTML}{D4D4D8}

\hypersetup{
    colorlinks=true,
    linkcolor=copifiemerald,
    urlcolor=copifisky,
    pdftitle={Brief stagiaire — Responsive Copifi},
    pdfauthor={Équipe Copifi},
    pdfsubject={Mission responsive mobile — SaaS Copifi}
}

\pagestyle{fancy}
\fancyhf{}
\fancyhead[L]{\small\textit{Brief stagiaire — Responsive Copifi}}
\fancyhead[R]{\small\thepage}
\renewcommand{\headrulewidth}{0.4pt}

\titleformat{\chapter}[display]
{\normalfont\huge\bfseries\color{copifiobsidian}}{\chaptertitlename\ \thechapter}{16pt}{\Huge}

\titleformat{\section}
{\normalfont\Large\bfseries\color{copifiemerald}}
{\thesection}{1em}{}

\titleformat{\subsection}
{\normalfont\large\bfseries\color{copifiobsidian}}
{\thesubsection}{1em}{}

\lstset{
    basicstyle=\ttfamily\small,
    backgroundcolor=\color{copificodebg},
    frame=single,
    rulecolor=\color{copificodeframe},
    breaklines=true,
    columns=fullflexible,
    keepspaces=true,
    showstringspaces=false,
    tabsize=2,
}

\newcommand{\filepath}[1]{\texttt{\small #1}}
\newcommand{\feat}[1]{\textbf{#1}}
\newcommand{\breakpoint}[2]{#1 & #2 \\}

\begin{document}

% ══════════════════════════════════════════════════════════════════
\begin{titlepage}
    \centering
    \vspace*{1.2cm}
    {\Huge\bfseries\color{copifiobsidian} Brief stagiaire\par}
    \vspace{0.6cm}
    {\LARGE\color{copifiemerald} Responsive — site Copifi\par}
    \vspace{0.4cm}
    {\large Mission jour 1 : rendre l'application utilisable sur mobile et tablette\par}
    \vspace{0.3cm}
    {\normalsize React + Inertia + Tailwind CSS — projet \filepath{saas\_cfo\_digital}\par}
    \vspace{1.4cm}
    \rule{0.78\textwidth}{1pt}
    \vspace{1cm}
    \begin{tabular}{>{\bfseries}r p{9.2cm}}
        Produit & Copifi (SaaS facturation + pilotage) \\
        Périmètre & Frontend uniquement (\filepath{resources/js/}) \\
        Stack & React 18, Inertia.js, Tailwind CSS 3, Vite \\
        Durée estimée & 1 journée (8 h) \\
        Branche Git cible & \filepath{feature/responsive-mobile} \\
        Version document & 1.0 — Juillet 2026 \\
    \end{tabular}
    \vfill
    \begin{minipage}{0.88\textwidth}
        \centering
        \small\itshape
        Ce document décrit l'objectif, la méthode de test, l'ordre de travail,
        les fichiers prioritaires et les livrables attendus en fin de journée.
        \textbf{Ne pas modifier} le backend PHP, les routes, ni la logique métier.
    \end{minipage}
\end{titlepage}

\tableofcontents
\cleardoublepage

% ============================================================
\chapter{Objectif de la mission}
% ============================================================

\section{Ce qu'on attend}

Rendre \textbf{tout le site Copifi} utilisable et lisible sur les viewports suivants :

\begin{itemize}[leftmargin=1.5em]
    \item \feat{320\,px} — iPhone SE (pire cas, prioritaire)
    \item \feat{375\,px} — iPhone standard
    \item \feat{768\,px} — tablette portrait
    \item \feat{1024\,px} — tablette paysage / bascule sidebar
    \item \feat{1280\,px+} — desktop (ne pas régresser l'existant)
\end{itemize}

\section{Definition of Done}

Une page est considérée \textbf{terminée} lorsque :

\begin{enumerate}[leftmargin=1.5em]
    \item Aucun scroll horizontal involontaire sur la page.
    \item Tous les boutons et liens sont cliquables au doigt (zone min.\ 44$\times$44\,px).
    \item Les textes ne débordent pas du viewport (pas de \texttt{overflow-x} caché).
    \item La navigation principale est accessible sans zoom.
    \item \texttt{npm run build} passe sans erreur après les modifications.
\end{enumerate}

% ============================================================
\chapter{Environnement et méthode}
% ============================================================

\section{Installation locale}

\begin{lstlisting}[language=bash]
cd saas_cfo_digital
composer install
npm ci
cp .env.example .env
php artisan key:generate
php artisan migrate
npm run dev          # terminal 1
php artisan serve    # terminal 2
\end{lstlisting}

\section{Outils de test obligatoires}

\begin{itemize}[leftmargin=1.5em]
    \item Chrome DevTools : \feat{F12} $\rightarrow$ \textbf{Toggle device toolbar} (Ctrl+Shift+M).
    \item Tester \textbf{chaque page modifiée} aux 5 largeurs listées au chapitre~1.
    \item Mode \textbf{Responsive} + mode \textbf{iPhone SE} explicite pour valider le pire cas.
\end{itemize}

\section{Breakpoints Tailwind (référence)}

Le projet utilise les breakpoints \textbf{par défaut} de Tailwind (aucun \texttt{screens} custom dans \filepath{tailwind.config.js}).

\begin{center}
\begin{tabular}{@{}lll@{}}
\toprule
\textbf{Préfixe} & \textbf{Min width} & \textbf{Usage typique} \\
\midrule
\breakpoint{(défaut)}{0 — mobile first}
\breakpoint{\texttt{sm:}}{640\,px}
\breakpoint{\texttt{md:}}{768\,px}
\breakpoint{\texttt{lg:}}{1024\,px — sidebar desktop, grilles larges}
\breakpoint{\texttt{xl:}}{1280\,px — dashboards larges}
\bottomrule
\end{tabular}
\end{center}

\textbf{Règle mobile first :} styles de base pour mobile, puis enrichissement avec \texttt{sm:}, \texttt{md:}, \texttt{lg:}.

% ============================================================
\chapter{Règles de développement}
% ============================================================

\section{À faire}

\begin{itemize}[leftmargin=1.5em]
    \item Utiliser \textbf{uniquement Tailwind CSS} (pas de lib CSS externe, pas de CSS inline sauvage).
    \item Réduire le padding mobile : \texttt{p-4 sm:p-6 lg:p-8} au lieu de \texttt{p-8} partout.
    \item Titres responsives : \texttt{text-2xl sm:text-3xl lg:text-4xl}.
    \item Toolbars filtres : \texttt{flex-col gap-3 sm:flex-row sm:items-center}.
    \item Tables : scroll horizontal explicite \textbf{ou} vue cartes en mobile (voir chapitre~\ref{ch:patterns}).
    \item Commits atomiques : \texttt{fix(responsive): factures index mobile}.
\end{itemize}

\section{Interdictions}

\begin{itemize}[leftmargin=1.5em]
    \item Modifier les contrôleurs PHP, les routes, les migrations, les services.
    \item Changer les couleurs, la typographie ou le logo (pas de refonte visuelle).
    \item Supprimer des colonnes de tableau sans alternative mobile.
    \item Utiliser \texttt{!important} ou des largeurs fixes en pixels sans breakpoint.
    \item Valider uniquement en zoom desktop 50\,\% (test device toolbar obligatoire).
\end{itemize}

% ============================================================
\chapter{Ordre de travail — priorités}
\label{ch:priorites}
% ============================================================

\section{Phase A — Layouts globaux (matin, $\sim$2\,h)}

Ces fichiers impactent \textbf{toutes} les pages de l'application.

\begin{longtable}{@{}p{5.2cm}p{4.2cm}p{5.6cm}@{}}
\toprule
\textbf{Fichier} & \textbf{Problème actuel} & \textbf{Action attendue} \\
\midrule
\endhead
\filepath{Layouts/AppDashboardLayout.jsx} &
Header chargé sur mobile ; \texttt{p-8} fixe &
\texttt{p-4 sm:p-6 lg:p-8} sur \texttt{<main>} ; simplifier header \texttt{< sm} \\
\filepath{Components/CfoPageShell.jsx} &
Double padding \texttt{-m-8 px-8} &
\texttt{px-4 sm:px-6 lg:px-8} \\
\filepath{Layouts/FacturationLayout.jsx} &
En-têtes page &
Vérifier \texttt{flex-col gap-4 sm:flex-row} \\
\filepath{Layouts/TestDashboardLayout.jsx} &
Sidebar fixe sans menu mobile &
Copier le drawer de \filepath{AppDashboardLayout.jsx} \\
\bottomrule
\end{longtable}

\textbf{Critère Phase A :} navigation OK à 320\,px, menu hamburger fonctionnel, contenu non coupé sous le header.

\section{Phase B — Landing ( $\sim$2\,h)}

\begin{longtable}{@{}p{5.5cm}p{8.8cm}@{}}
\toprule
\textbf{Fichier} & \textbf{Action attendue} \\
\midrule
\endhead
\filepath{Pages/Welcome.jsx} &
Menu hamburger \texttt{< md} avec ancres ; hero \texttt{text-3xl sm:text-4xl lg:text-5xl} ; KPIs \texttt{grid-cols-1 sm:grid-cols-3} ; atténuer mockup 3D \texttt{< md} \\
\filepath{Components/Landing/LandingChatWidget.jsx} &
\texttt{min-h-[320px] sm:min-h-[420px] lg:min-h-[520px]} \\
\bottomrule
\end{longtable}

\section{Phase C — Pages listes / tableaux ($\sim$3\,h)}

\begin{longtable}{@{}p{5.8cm}p{2.2cm}p{6.3cm}@{}}
\toprule
\textbf{Fichier} & \textbf{min-w table} & \textbf{Priorité} \\
\midrule
\endhead
\filepath{Pages/FacturesFournisseurs/Index.jsx} & 960\,px & \textcolor{copifiemerald}{\textbf{Haute — 0 responsive actuel}} \\
\filepath{Pages/Factures/Index.jsx} & 980\,px & Haute \\
\filepath{Pages/Devis/Index.jsx} & 960\,px & Haute \\
\filepath{Pages/FinFlow/Payments/Index.jsx} & 900\,px & Moyenne \\
\filepath{Pages/FinFlow/Catalogue/Index.jsx} & 900\,px & Moyenne \\
\filepath{Pages/FinFlow/Clients/Index.jsx} & 760\,px & Moyenne \\
\filepath{Components/FinFlow/DashboardHome.jsx} & 640\,px & Moyenne \\
\bottomrule
\end{longtable}

\textbf{Option rapide :} conserver le scroll + ajouter un conteneur \texttt{overflow-x-auto}.\\
\textbf{Option recommandée :} vue \textbf{cartes} \texttt{< md}, table \texttt{hidden md:block}.

\section{Phase D — Formulaires lourds ($\sim$3\,h)}

\begin{itemize}[leftmargin=1.5em]
    \item \filepath{Pages/FinFlow/InvoiceCreate.jsx} — lignes en cartes \texttt{< lg} ; barre d'actions empilée ou menu \texttt{⋮} ; preview sous le formulaire.
    \item \filepath{Pages/Devis/Create.jsx} — même traitement que la facture.
\end{itemize}

\textbf{Critère :} créer une facture/devis entièrement sur iPhone 375\,px sans zoom navigateur.

\section{Phase E — Dashboards et reste ($\sim$2\,h)}

\begin{itemize}[leftmargin=1.5em]
    \item \filepath{Pages/Dashboard.jsx} — KPIs \texttt{grid-cols-1 sm:grid-cols-2 xl:grid-cols-4}
    \item \filepath{Pages/Facturation/Dashboard.jsx} — via \filepath{DashboardHome.jsx}
    \item \filepath{Pages/FinFlow/Settings/Index.jsx} — sections empilées
    \item \filepath{Pages/Copilot.jsx} — stack vertical \texttt{< lg}
    \item \filepath{Pages/Auth/AuthShell.jsx} — vérifier 320\,px
    \item \filepath{Pages/Admin/Dashboard.jsx} — table scroll ou cartes
\end{itemize}

% ============================================================
\chapter{Patterns réutilisables}
\label{ch:patterns}
% ============================================================

\section{Table $\rightarrow$ cartes mobile}

\begin{lstlisting}[language=Java]
{/* Mobile : cartes */}
<div className="space-y-3 md:hidden">
  {items.map((item) => (
    <div key={item.id}
         className="rounded-xl border border-white/10 bg-white/5 p-4">
      <p className="font-medium text-white">{item.reference}</p>
      <p className="text-sm text-zinc-400">{item.client}</p>
      <p className="mt-2 text-lg font-semibold text-emerald-400">
        {item.total}
      </p>
    </div>
  ))}
</div>

{/* Desktop : table */}
<div className="hidden md:block overflow-x-auto">
  <table className="min-w-full ...">...</table>
</div>
\end{lstlisting}

\section{Toolbar filtres}

\begin{lstlisting}[language=Java]
<div className="flex flex-col gap-3 sm:flex-row sm:items-center sm:justify-between">
  <input className="w-full sm:max-w-xs" />
  <div className="flex flex-wrap gap-2">{/* boutons */}</div>
</div>
\end{lstlisting}

\section{Padding page cohérent}

\begin{lstlisting}[language=Java]
<main className="flex-1 p-4 sm:p-6 lg:p-8">
\end{lstlisting}

% ============================================================
\chapter{Inventaire des fichiers}
% ============================================================

\section{Layouts (\filepath{resources/js/Layouts/})}

\begin{itemize}[leftmargin=1.5em]
    \item \filepath{AppDashboardLayout.jsx} — shell principal (sidebar + header)
    \item \filepath{FacturationLayout.jsx} — wrapper facturation
    \item \filepath{GuestLayout.jsx} — auth legacy
    \item \filepath{TestDashboardLayout.jsx} — prototype (à corriger)
\end{itemize}

\section{Pages (\filepath{resources/js/Pages/}) — 31 fichiers}

\textbf{Racine :} \filepath{Welcome.jsx}, \filepath{Dashboard.jsx}, \filepath{Copilot.jsx}, \filepath{FinancialEntry.jsx}, pages Test*.\\
\textbf{Auth :} \filepath{Auth/AuthShell.jsx}, Login, Register, etc.\\
\textbf{Facturation :} \filepath{Facturation/Dashboard.jsx}\\
\textbf{FinFlow :} InvoiceCreate, Clients, Catalogue, Payments, Settings, Placeholder\\
\textbf{Documents :} Devis (Index, Create), Factures (Index), FacturesFournisseurs (Index)\\
\textbf{Admin / Profile :} \filepath{Admin/Dashboard.jsx}, \filepath{Profile/Edit.jsx}

\section{Composants partagés}

\begin{itemize}[leftmargin=1.5em]
    \item \filepath{Components/FinFlow/} — KPI, modales, dashboard home (15 fichiers)
    \item \filepath{Components/CfoPageShell.jsx}, \filepath{CfoPageBackground.jsx}
    \item \filepath{Components/Landing/LandingChatWidget.jsx}
    \item \filepath{Components/Dashboard/DashboardChatWidget.jsx}
    \item \filepath{Components/ui/} — login, logo
\end{itemize}

% ============================================================
\chapter{Checklist de validation}
% ============================================================

\section{Checklist stagiaire (à cocher en fin de journée)}

\begin{itemize}[leftmargin=1.5em]
    \item[$\square$] \filepath{Welcome.jsx} — 320, 375, 768, 1024
    \item[$\square$] \filepath{AppDashboardLayout.jsx} — menu + header
    \item[$\square$] \filepath{Dashboard.jsx} (pilotage)
    \item[$\square$] \filepath{Facturation/Dashboard.jsx}
    \item[$\square$] \filepath{Factures/Index.jsx}
    \item[$\square$] \filepath{Devis/Index.jsx} + \filepath{Devis/Create.jsx}
    \item[$\square$] \filepath{FinFlow/InvoiceCreate.jsx}
    \item[$\square$] \filepath{FinFlow/Clients/Index.jsx}
    \item[$\square$] \filepath{FinFlow/Catalogue/Index.jsx}
    \item[$\square$] \filepath{FinFlow/Payments/Index.jsx}
    \item[$\square$] \filepath{FinFlow/Settings/Index.jsx}
    \item[$\square$] \filepath{FacturesFournisseurs/Index.jsx}
    \item[$\square$] \filepath{Copilot.jsx}
    \item[$\square$] Auth (Login, Register)
    \item[$\square$] \filepath{Profile/Edit.jsx}
\end{itemize}

\section{Validation tuteur (5 min)}

\begin{enumerate}[leftmargin=1.5em]
    \item iPhone 375\,px : landing + login + liste factures + création facture.
    \item Aucun scroll horizontal involontaire.
    \item Boutons cliquables au doigt.
    \item \texttt{npm run build} OK.
\end{enumerate}

% ============================================================
\chapter{Livrables fin de journée}
% ============================================================

\begin{enumerate}[leftmargin=1.5em]
    \item \textbf{Branche Git} : \filepath{feature/responsive-mobile} (PR ou merge request).
    \item \textbf{Screenshots} avant/après pour 3 pages :
    \begin{itemize}
        \item Welcome (landing)
        \item Factures/Index
        \item InvoiceCreate
    \end{itemize}
    à 375\,px de large.
    \item \textbf{Liste des fichiers modifiés} (commentaire PR ou fichier \filepath{RESPONSIVE-NOTES.md}).
    \item \textbf{Notes} : pages non terminées + blocages éventuels.
\end{enumerate}

\vfill
\begin{center}
\small\color{copifizinc500}
Document interne Copifi — Brief stagiaire responsive — Juillet 2026\\
Compilable avec \texttt{pdflatex} (2 passes pour la table des matières).
\end{center}

\end{document}
